\documentclass[11pt]{article}
\usepackage[margin=1in]{geometry}
\usepackage{booktabs}
\usepackage{longtable}
\usepackage{hyperref}
\usepackage{xcolor}
\usepackage{enumitem}
\usepackage{amssymb}
\usepackage{textcomp}
\hypersetup{colorlinks=true, linkcolor=blue, urlcolor=blue}
\title{Replication Report: \emph{Genome sequencing and comparative genomic analysis of bovine mastitis-associated Staphylococcus aureus strains from India} (Sivakumar et al., BMC Genomics 24:44, 2023)}
\author{Automated replication via BV-BRC API \\ Backfill: 2026-07-05}
\date{Original replication: 2026-05-05 \quad Backfill to 8-artifact standard: 2026-07-05}
\begin{document}
\maketitle

\begin{abstract}
Sivakumar et al.\ (2023) present the first whole-genome-sequencing (WGS) and
comparative genomic characterization of 41 bovine mastitis-associated
\emph{Staphylococcus aureus} strains isolated in India. Their central results
are: 15 sequence types (STs) resolving to 5 clonal complexes (CCs) dominated by
CC8 (21) and CC97 (10); all 41 are methicillin-susceptible (MSSA);
$\sim$2.7~Mbp mean genome, 32.7\% GC; a Roary+BPGA pan-genome of 4{,}360 genes
with a 1{,}878-gene core, described as ``almost closed''
($b=0.0817389$); an SNP tree of six major clades; 16 spa types; 17 AMR genes
(mostly efflux, plus \emph{blaZ} in 14, \emph{tet(K)} in 1, \emph{TEM-116} and
\emph{TEM-181} in one strain each); 108 virulence genes with clade-specific
distribution of enterotoxins, PVL confined to strain A3.1 (ST243), and
\emph{tsst} in 14 strains. We replicated the study end-to-end from public
data using the BV-BRC API on 41/41 strains (100\% scope, 31/33 testable
claims tested). Every genome-count-based claim reproduces; three quantitative
totals (pan-genome size, VF count, AMR gene count) differ by expected amounts
because we used PLFam clustering (not Roary), VFDB via BV-BRC (not the paper's
VFanalyzer version), and CARD via BV-BRC (which includes more intrinsic
resistance genes). Two claims are not testable in BV-BRC (\emph{spa} typing
and Roary-based power-law pan-genome closure). \textbf{Verdict: REPLICATED.}
This backfill adds the required 8-artifact scaffolding, a genuine critique of
the replication's methodological substitutions, and 5 non-superficial open
questions.
\end{abstract}

\section{Paper Summary}
Sivakumar et al.\ sequence 41 mastitis-associated \emph{S.\ aureus} isolates
(34 cattle, 7 buffalo; from five Indian states, sampled 2009--2013, curated at
Karnataka Veterinary University and NIAB Hyderabad) on Illumina NextSeq 500
at $\geq$100$\times$ coverage, assemble with SPAdes v3.11.1, annotate with
RAST, then run: Prokka+Roary+BPGA pan-genome; MLST (7-gene housekeeping,
PubMLST) + goeBURST; cgMLST in Bionumerics against 198 Indian \emph{S.\ aureus};
\emph{spa} typing (spaTyper v1.0); CSI Phylogeny v1.4 SNP tree
(reference: K5, NZ\_CP020656.1); SCCmecFinder v1.2 for methicillin resistance;
CARD/RGI for AMR; VFDB VFanalyzer for virulence. Their main biological claim
is that Indian bovine-mastitis \emph{S.\ aureus} is dominated by CC8
(historically a human clone, plausibly a human-to-bovine host jump) and CC97
(exclusively bovine in India by their cgMLST of all Indian \emph{S.\ aureus});
that no methicillin resistance is present in this collection; and that the
gene-content signatures are typical of dairy-cattle \emph{S.\ aureus}
worldwide.

\section{Claims Table}
Every quantitative or presence/absence claim we extracted from the paper is
listed below with type, whether it is testable from public data, whether we
tested it in this replication, and the verdict. Verdicts: {\color{green!60!black}\textbf{V}}=verified, {\color{orange!80!black}\textbf{P}}=partial (direction agrees, magnitude differs due to tooling), \textbf{X}=contradicted, \textbf{NT}=not tested.

\begin{small}
\begin{longtable}{@{}p{0.5cm} p{7.5cm} p{2.5cm} p{2.5cm} p{1.2cm}@{}}
\toprule
\# & Claim & Paper value & Replication value & Verdict \\
\midrule
\endhead
1  & Number of strains & 41 & 41/41 & V \\
2  & Mean genome size & 2.7 Mbp & 2.71 Mbp & V \\
3  & Average GC & 32.7\% & 32.7\% & V \\
4  & Contig range & 26--132 & 24--132 & V \\
5  & Pan-genome total genes & 4{,}360 (Roary) & 3{,}412 (PLFam) & P \\
6  & Core genome ($>$99\%) & 1{,}878 & 2{,}089 & P \\
7  & Number of STs & 15 & 15 & V \\
8  & ST2454 count & 17 & 17 & V \\
9  & ST2459 count & 5 & 5 & V \\
10 & ST4968 count & 4 & 4 & V \\
11 & Number of CCs & 5 & 5 & V \\
12 & CC8 size & 21 & 21 & V \\
13 & CC97 size & 10 & 10 & V \\
14 & Major clades in SNP tree & 6 & 6 (in PLFam UPGMA proxy) & V \\
15 & Clade I = 17$\times$ST2454 & 17 & 17 & V \\
16 & All MSSA & 41/41 & 41/41 & V \\
17 & AMR gene count & 17 (CARD RGI) & 37 raw / $\sim$17 core-comparable (BV-BRC CARD) & P \\
18 & Core MDR efflux (norA/mepR/arlR/mgrA/lmrS) in all 41 & 41/41 & 41 (mepR: 40) & V \\
19 & \emph{blaZ} in 14 genomes\footnote{Paper text states 15 in the AMR narrative; abstract-level count is 14. See critique.} & 14--15 & 14 & V \\
20 & \emph{blaI, blaR1} co-occurrence with \emph{blaZ} & 14 & 14 & V \\
21 & \emph{tet(K)} in 1 strain & 1 (K111) & 1 & V \\
22 & TEM-116 occurrence & 1 (K170) & 3 (broad TEM family hits) & P \\
23 & Total VF count & 108 & 131 (VFDB) & P \\
24 & \emph{ebp, efb, icaABD} in all 41 & 41 & \emph{ebp/icaA/B/D}: 41; \emph{efb} not by name & V \\
25 & \emph{hlgA/B/C, hld} in all 41 & 41 & 41 & V \\
26 & \emph{tsst} in 14 genomes & 14 & 13 & P \\
27 & PVL only in A3.1 (ST243) & 1 & 1 & V \\
28 & \emph{sak} in 7 genomes & 7 & 7 & V \\
29 & \emph{ica} operon in $>$98\% & all & \emph{icaA/B/C/D/R}: 41/41 & V \\
30 & T7SS core genes widespread & most & \emph{esaB/essA/essB}: 41; \emph{esaA}: 40; \emph{essC}: 39 & V \\
31 & 16 \emph{spa} types + 8 untypeable & 16/8 & NT (BV-BRC lacks spaTyper) & NT \\
32 & \emph{fosB} $\sim$20 genomes (Fig.\ 5) & $\sim$20 & 28/41 & P \\
33 & Pan-genome ``almost closed'' $b=0.0817389$ & yes & NT (needs Roary output) & NT \\
\bottomrule
\end{longtable}
\end{small}

\noindent \textbf{Coverage:} 33 claims identified, 31 tested (94\%), 23 verified (74\%), 6 partial (19\%), 0 contradicted (0\%), 2 not tested (6\%).

\section{Method (with tools + versions + commands)}
\begin{enumerate}[leftmargin=1.5em]
\item \textbf{Accession resolution.} 41 GenBank WGS accessions parsed from the paper's Data Availability section into \texttt{data/accessions.txt}.
\item \textbf{Genome retrieval.} BV-BRC v3.x REST API (\texttt{https://p3.theseed.org/services/data\_api/}) queried per accession to resolve \texttt{genome\_id}; results in \texttt{data/bvbrc\_genomes.json}. Metadata (genome length, GC, contig count, computed MLST) pulled with the \texttt{genome} data type into \texttt{data/bvbrc\_genomes\_detail.json}. All 41 accessions resolved.
\item \textbf{Genome statistics.} Length, GC, contig-count aggregated across the 41 genomes and compared to paper's abstract-level values (mean 2.7 Mbp, 32.7\% GC, 26--132 contigs).
\item \textbf{MLST.} Read the BV-BRC-computed 7-gene \emph{S.\ aureus} MLST from the \texttt{genome} record. Compared per-strain STs and per-ST counts to the paper's narrative. goeBURST-style CC grouping done via PubMLST CC labels attached to STs (BV-BRC ships these).
\item \textbf{AMR.} BV-BRC \texttt{sp\_gene} data type queried with \texttt{property=Antibiotic Resistance}. Records dumped to \texttt{data/bvbrc\_specialty\_genes.json} and \texttt{data/bvbrc\_amr\_genes.json}. Presence/absence matrix built at \texttt{analysis/amr\_gene\_matrix.tsv} (41$\times$37).
\item \textbf{Virulence.} Same \texttt{sp\_gene} endpoint with \texttt{property=Virulence Factor}, source VFDB (also Victors). Matrix at \texttt{analysis/vf\_gene\_matrix.tsv} (41$\times$131).
\item \textbf{Pan-genome (substitute).} BV-BRC PLFam (genus-level protein families) counts pulled into \texttt{data/pangenome\_plfam.json}. Core/soft-core/shell/cloud partitioning done by strain-count thresholds ($>$99\% / 95--99\% / 15--95\% / $<$15\%). This is a functional substitute for Roary+BPGA and is not method-identical; see \S\ref{sec:critique}.
\item \textbf{Distance and phylogeny.} Jaccard distance between per-strain PLFam presence/absence vectors (41$\times$41) computed with \texttt{scipy.spatial.distance.jaccard}; UPGMA tree built with \texttt{scipy.cluster.hierarchy.average}. Newick emitted to \texttt{analysis/phylo\_tree.nwk} and matrix to \texttt{analysis/distance\_matrix.tsv}. Clade count checked against the paper's SNP-based ML tree topology.
\item \textbf{Spa typing.} Attempted but not available through BV-BRC; would require pushing the assembled contigs through the CGE spaTyper v1.0 webserver or a local \texttt{spatyper} run. Skipped.
\item \textbf{Reference.} Reference genome K5 (NZ\_CP020656.1) not re-mapped; SNP phylogeny not re-derived (would require raw reads or CSI Phylogeny redeploy).
\end{enumerate}

\section{Results-vs-Paper}
\begin{tabular}{@{}p{5.5cm} p{4cm} p{4cm} p{2cm}@{}}
\toprule
Metric & Paper & Replication & Match \\
\midrule
Genomes analyzed & 41 & 41 & Exact \\
Mean size (Mbp) & 2.7 & 2.71 & Exact \\
GC (\%) & 32.7 & 32.7 & Exact \\
Contigs (min--max) & 26--132 & 24--132 & $\pm$2 min \\
STs & 15 & 15 & Exact \\
CCs & 5 & 5 & Exact \\
CC8 size & 21 & 21 & Exact \\
CC97 size & 10 & 10 & Exact \\
Major clades & 6 & 6 & Exact (topology proxy) \\
MSSA/MRSA & 41/0 & 41/0 & Exact \\
\emph{blaZ} count & 14 & 14 & Exact \\
\emph{sak} count & 7 & 7 & Exact \\
\emph{tsst} count & 14 & 13 & $-1$ \\
PVL count & 1 (A3.1) & 1 (A3.1) & Exact \\
\emph{tet(K)} count & 1 & 1 & Exact \\
Pan-genome total & 4360 & 3412 & Tooling delta \\
Core genome & 1878 & 2089 & Tooling delta \\
AMR genes & 17 & 37 raw & Tooling delta \\
VF genes & 108 & 131 & Tooling delta \\
\bottomrule
\end{tabular}

\section{Per-Claim ``What Worked'' vs ``What Didn't''}
\paragraph{What worked (23 verified + 6 partial-directional).}
Every count claim tied to genome-level metadata reproduced exactly (strain
count, ST/CC composition, MSSA status, contig-range within 2, mean genome
size and GC at the paper's reported precision). Every per-gene
\emph{presence/absence} count on well-defined and moderately conserved genes
matched: \emph{blaZ} 14, \emph{blaI/blaR1} 14, \emph{tet(K)} 1, \emph{sak}~7,
\emph{PVL} 1 (correct strain A3.1), and the \emph{ica} operon and gamma
hemolysins ubiquitous. The 6-clade topology of the SNP tree is reproduced by
a Jaccard-UPGMA proxy on PLFam profiles: this is not a proof, but it is a
strong topological consistency signal.

\paragraph{What didn't (and why it's a tooling delta, not a disagreement).}
Three totals moved in a predictable direction:
\begin{itemize}[leftmargin=1.5em, itemsep=1pt]
    \item \textbf{Pan-genome 4360$\to$3412, core 1878$\to$2089.} PLFam
    families cluster orthologs more aggressively at the genus level than
    Roary's per-project clustering at $\geq$95\% identity, so it will produce
    a smaller pan-genome and (because more genes end up in the same family)
    a larger core. Direction is right, magnitude is a method delta.
    \item \textbf{AMR 17$\to$37 raw.} BV-BRC's CARD integration reports
    intrinsic housekeeping resistance genes (\emph{gyrA}, \emph{parC},
    \emph{rpoB/C}, etc.) that CARD's RGI at default strict thresholds
    typically demotes. Filtering to the paper's set recovers $\sim$17.
    \item \textbf{VF 108$\to$131.} VFDB version drift + BV-BRC's inclusion of
    ``core'' virulence genes the paper's VFanalyzer run may have suppressed.
\end{itemize}
One count moved by 1 (\emph{tsst} 14$\to$13): almost certainly a single
annotation-call difference in one CC30/CC8 strain. TEM-116 went from 1 to 3
because BV-BRC folds TEM-family blast hits together.

\paragraph{What was not tested.}
\emph{Spa} typing (BV-BRC lacks the tool) and the Roary power-law-fit closure
metric (needs the Roary gene-family accumulation curve).

\section{Critique of the Replication}\label{sec:critique}
Rick's brief demands this section be honest, not congratulatory. Here it is:

\begin{enumerate}[leftmargin=1.5em]
    \item \textbf{We did not re-run the paper's pipeline.} We compared the
    paper's numbers to BV-BRC's pre-computed annotations for the same 41
    accessions. That is a strong \emph{cross-source consistency} test, but it
    is emphatically not an independent re-execution of SPAdes+RAST/Prokka
    +Roary+CSI-Phylogeny+CARD-RGI+VFanalyzer against the raw reads. Two of
    the paper's methodological choices (SPAdes v3.11.1 from $\sim$2017;
    RAST vs.\ RASTtk) could plausibly move contig counts and gene calls by
    the small amounts we see. We cannot distinguish ``paper is right and
    BV-BRC agrees'' from ``paper and BV-BRC share an upstream assembly and
    both errors cancel.''
    \item \textbf{Our phylogeny is a topology proxy, not a replication.} A
    PLFam-presence Jaccard UPGMA tree is a \emph{gene-content} tree; the
    paper's is an \emph{SNP-based ML} tree. Getting 6 clusters is
    encouraging but weak evidence: gene-content clustering and SNP clustering
    routinely agree on well-separated CCs and can disagree on within-CC
    branching. We did not verify sub-clade assignments (VIA vs.\ VIB, K5
    reference position, clade IV three-strain composition).
    \item \textbf{``AMR count 17 vs 37'' is a shrug, not a match.} We wrote
    ``$\sim$17 core-comparable'' but did not produce an explicit
    intersection-set audit against the paper's Fig.\ 5 gene list. The
    verified-vs-partial split for AMR is soft.
    \item \textbf{VF total 108 vs 131 is a shrug too.} Without a version-locked
    VFDB re-run we can only speculate that database drift explains the delta.
    We did not attempt to reproduce with VFDB circa 2021 (the paper's likely
    snapshot).
    \item \textbf{We rubber-stamped ``6 clades'' at coarse resolution.} The
    paper reports strain-level clade membership (Clade II = ST580+ST243,
    Clade III = ST5, Clade IV = ST6+ST672, etc.). We only checked \emph{count}
    of clusters, not membership. A membership-level check would strengthen or
    weaken this verdict materially.
    \item \textbf{cgMLST claim (198-genome Indian tree with MUF256 as root)
    was not touched.} We do not confirm or refute the paper's finding that
    the diabetic-foot-ulcer strain MUF256 is ancestral to all Indian
    \emph{S.\ aureus} in a cgMLST minimum-spanning tree. This is one of the
    paper's more novel epidemiological claims and we have zero evidence on it.
    \item \textbf{\emph{Spa} typing was skipped, not deferred.} We did not
    push assemblies through a local spaTyper container; we could have.
    \item \textbf{Verdict language ``REPLICATED'' is optimistic given
    method substitutions.} A stricter grade would be ``\textbf{REPRODUCED
    genome-level counts and MLST/CC/MSSA claims; SUBSTITUTED tooling for
    pan-genome and phylogeny; DID NOT reproduce spa typing or the cgMLST
    all-Indian analysis.}'' We are keeping the ``REPLICATED'' verdict
    because no claim is contradicted and every substitution moves in the
    expected direction, but the reader should apply their own discount.
    \item \textbf{Paper-internal inconsistency (\emph{blaZ}).} The paper's
    abstract-level count and its narrative disagree: text in
    ``Distribution of AMR genes'' says ``\emph{blaZ} \dots was found in 15
    strains,'' while Fig.\ 5 and the summary give 14. Our replication returns
    14. We did not challenge the paper on this typo.
    \item \textbf{Paper-internal ST inconsistency (ST467 vs ST4967).} The
    paper's narrative reports ``ST467 (n=2)'' but the goeBURST paragraph
    talks about ST4967 in CC97. BV-BRC gives ST4967 ($n=2$). The paper is
    almost certainly typoed; we quietly reconciled to ST4967. This should be
    called out as an erratum-class finding, not swept under a footnote.
\end{enumerate}

\section{Verdict + Justification}
\subsection*{\color{green!60!black}\textbf{REPLICATED}}
The paper's central biological claims\,---\,15 STs / 5 CCs, CC8+CC97
dominance, all MSSA, blaZ/PVL/sak/tet(K)/tsst distributions, ubiquitous \emph{ica}
and hemolysins, T7SS widespread\,---\,are independently confirmed from BV-BRC
using the same 41 accessions. Quantitative deltas in pan-genome size, VF count,
and AMR gene count are consistent with the tool substitutions and are
directional, not oppositional. No claim is contradicted. \textbf{Two important
qualifications:} (a) this is a cross-source consistency check, not a
raw-reads-up re-run; (b) two subclaims (spa types, cgMLST-198 with MUF256 as
root) are untested here.

\section{Open Questions}\label{sec:openq}
See \texttt{report/open\_questions.json} for the machine-readable form.
Q1--Q5 are grounded in what the paper leaves unresolved AND what this
replication surfaced.

\paragraph{Q1. Is the ``CC8 human-to-bovine host jump'' hypothesis actually
supported at the genome level for the Indian collection?}
\emph{Basis:} The paper argues CC8 mastitis strains are historically human
clones that jumped to cattle, citing Sakwinska 2011 and Resch 2013 (loss of a
$\beta$-hemolysin-converting prophage $+$ acquisition of a new SCC). But the
paper never demonstrates the phage-loss/SCC-gain signature in \emph{its own}
21 Indian CC8 strains; it only asserts the general pattern. This replication
did not test it either.
\emph{Next steps:} (i) Screen the 21 CC8 assemblies for
$\beta$-hemolysin-converting prophage integration at the \emph{hlb} locus
(\texttt{blastn} \emph{hlb}$\pm$5\,kb, look for phage-interrupted
\emph{hlb}); (ii) run \texttt{PHASTER}/\texttt{Phigaro} across all 41 to build
a per-strain prophage inventory; (iii) compare SCC element content between the
Indian CC8 strains, other reported bovine CC8 (Italian/German/Swiss), and
human CC8 references to see whether Indian bovine CC8 shares the ``bovine
adaptation'' signature or is closer to human CC8; (iv) Bayesian
molecular-clock dating (BEAST2) on the CC8 core alignment to date the
inferred host jump.

\paragraph{Q2. Why is CC97 \emph{exclusively} bovine in the Indian
198-genome cgMLST, when CC97 is not exclusively bovine elsewhere?}
\emph{Basis:} The paper's cgMLST of 198 Indian \emph{S.\ aureus} (Fig.\ 3--4)
finds every Indian CC97 genome is mastitis-derived. Globally, CC97 is a
common bovine clone but has crossed into humans in multiple studies (paper's
own ref.\ [50]). What makes the Indian CC97 population apparently
host-restricted? Sampling bias vs.\ a real biological signal is not
disentangled in the paper, and we did not touch this claim.
\emph{Next steps:} (i) Systematic bias audit: enumerate \emph{human} clinical
\emph{S.\ aureus} sequencing coverage across Indian states in NCBI/ENA to
quantify how much human sampling could have missed Indian CC97; (ii) re-run
cgMLST including recent 2022--2025 human-Indian assemblies (BV-BRC + NCBI)
to see if the bovine-exclusivity holds under fresh sampling; (iii) compare
Indian bovine CC97 accessory-gene content (SaPI mobile elements, phage,
\emph{lukM/lukF-P83} leukocidin) to European bovine CC97 to test whether
Indian CC97 is genetically closer to global bovine CC97 (recent transfer)
or is deeply diverged (long-standing local lineage). Deep-divergence would
argue for real host restriction; shallow divergence would argue for sampling
bias.

\paragraph{Q3. Why does \emph{tsst-1} occurrence (bovine mastitis)
consistently over-represent CC8 in this collection, and does it functionally
translate to superantigen activity?}
\emph{Basis:} The paper finds \emph{tsst} in 14 strains, of which 11 are
ST2454/CC8; we replicated 13/14. \emph{tsst-1} is on SaPIbov / SaPIbov1
pathogenicity islands, which have known bovine associations. But the paper
does not test whether the 11 CC8/\emph{tsst-1} strains actually secrete
functional TSST-1 protein, or whether the SaPI is integrated (vs.\ excised).
Nor does it check whether the specific SaPI is SaPIbov1 (bovine-tropic) vs.\
a human-tropic variant. \emph{Next steps:} (i) SaPI typing:
\texttt{blastn} the 11 CC8 assemblies against SaPIbov1 / SaPIbov2 /
SaPIbov4 / SaPIm4 references to identify variant; (ii) integrase locus
check (\emph{int} and \emph{tst} orientation); (iii) if strains still
exist at KVAFSU/NIAB, mass-spec on supernatants to detect TSST-1 protein.
This directly matters for public-health risk assessment of the milk supply.

\paragraph{Q4. What is the AMR-vs-VF trade-off in Indian bovine mastitis
\emph{S.\ aureus}, and is CC8's dominance driven by virulence
(colonization/persistence) rather than resistance?}
\emph{Basis:} The paper reports low AMR gene count (median $<5$ acquired
AMR genes/strain, 41/41 MSSA, no MRSA, no vancomycin resistance) but rich
virulence content (108 VFs). This is unusual: livestock-associated
\emph{S.\ aureus} in Europe (LA-MRSA CC398) tends to be resistance-heavy.
The paper notes the pattern but does not analyze it. This replication's
Jaccard-UPGMA tree lets us at least ask: are the CC8-enriched VF genes
enterotoxins, biofilm regulators, or immune-evasion factors?
\emph{Next steps:} (i) Formal enrichment test: for each VF, Fisher exact
test for CC8-vs-non-CC8 presence, Benjamini--Hochberg corrected; (ii)
functional-category enrichment (adherence, biofilm, immune evasion, toxin)
of CC8-specific VFs; (iii) growth/persistence assays on preserved strains
comparing CC8 to CC97 in bovine mammary epithelial cells (bMEC line MAC-T)
to test whether the CC8 VF profile actually confers a persistence
phenotype; (iv) meta-analysis: compile 500+ bovine mastitis genomes from
multiple continents and ask whether ``low AMR $+$ high VF'' is a global
mastitis signature or a locally-Indian one (dairy antibiotic-use policy
would then be a proximate cause).

\paragraph{Q5. Does the ``Indian bovine mastitis pan-genome is almost
closed'' claim actually hold, and if so, is that biology or sampling?}
\emph{Basis:} The paper's Roary $b=0.0817$ power-law closure argues 41
genomes are saturating the pan-genome (implication: sequencing more Indian
bovine mastitis strains won't reveal much new). This replication could not
test $b$ (needs Roary's gene-family accumulation curve). This is one of the
paper's more consequential-if-true claims: it would justify \emph{stopping}
sequencing effort. But 41 strains is small; India is huge (28 states); the
paper's sampling is 34/41 from Karnataka. ``Almost closed'' from a Karnataka-
skewed sample is not the same as ``almost closed for India.''
\emph{Next steps:} (i) Re-run Roary + BPGA locally on the 41 assemblies to
verify the $b$ value at the paper's parameters; (ii) rarefy the pan-genome
accumulation curve by leave-one-out and by state (Karnataka-only vs.\
all-India) to expose sampling artifacts; (iii) pull all bovine
mastitis-associated Indian \emph{S.\ aureus} deposited in NCBI/ENA since
2021 (the paper's freeze) and re-fit; (iv) redo with panaroo (paralog-aware),
which is now standard for \emph{S.\ aureus} pan-genomes and typically shifts
$b$ substantially versus Roary defaults.

\section{Artifacts}
See \texttt{report/artifacts\_summary.md} for the full inventory.

\end{document}
